\chapter{Discussion and Analysis}
\label{ch:evaluation}

This chapter interprets and analyses the results presented in Chapter~\ref{ch:results}. It explains what the findings mean in the context of the research objectives, situates them within the existing multi-agent literature, evaluates their significance, and identifies the main limitations of this work.

% ================================================================
\section{Interpretation of Results}
\label{sec:interpretation}
% ================================================================

\subsection{Cooperative Multi-Agent Coordination via PPO}
\label{subsec:disc_ppo}

Objective O3 is fully achieved. The MuJoCo Soccer 2v2 environment was successfully configured with our custom shared-policy PPO architecture. The evaluation results in Figure~\ref{fig:soccer_eval} demonstrate that the agents accumulate non-trivial shaped rewards across episodes, confirming that they have learned to coordinate, pursue the ball, and orient themselves towards the goal mouth.

The training results confirm that policy parameter sharing via SuperSuit effectively quadrupled the training data per update step, as all four agents' observations contribute to the same PPO gradient. This is highly consistent with findings in multi-agent literature~\citep{gupta2017cooperative}, which demonstrate that parameter sharing in symmetric multi-agent tasks significantly accelerates convergence by reducing the dimensional search space.

\subsection{Effectiveness of Custom Reward Shaping}
\label{subsec:disc_reward_shaping}

The custom reward shaping wrapper (Eq.~\eqref{eq:total_reward}) proved absolutely critical. In the sparse-reward baseline, the agents fail to register any learning trajectory due to the extreme rarity of goal-scoring events. The distance-to-ball penalty forces the agents to stay close to the ball; the velocity reward reinforces purposeful ball-chasing behavior; and the goal velocity reward creates an immediate incentive to advance the ball towards the opponent's goal. This dense reward structure closely follows the reward shaping principles discussed in literature~\citep{ng1999policy}, transforming a sparse exploration problem into a continuous directional gradient.

\subsection{Spatial Coverage Heatmaps and Teammate Spacing Dynamics}
\label{subsec:disc_spatial_analytics}

Objective O4 is fully achieved. The 2D pitch spatial density heatmaps (Figure~\ref{fig:soccer_heatmap}) provide strong qualitative proof of tactical learning:
\begin{itemize}
    \item \textbf{Territorial Division}: Team Red (Attackers) spent the majority of their time in the opposition's half, peaking around $[-4.0, -1.5]$ and $[-2.0, 2.0]$. Conversely, Team Blue (Defenders) maintained dense positioning in their own half (around coordinates $[4.0, 1.5]$ and $[2.0, -2.0]$). This highlights that PPO successfully discovered spatial divisions, matching classical offensive and defensive roles.
    \item \textbf{Optimal Teammate Spacing}: The average team spread was measured at **5.39m** for Team Red and **5.47m** for Team Blue. This spacing is highly optimal: it is large enough to prevent teammate collisions (clustering) but small enough to maintain close cooperative support.
    \item \textbf{Possession Balance}: The balanced possession (Red 48.6\% vs Blue 51.4\%) indicates highly stable, competitive match dynamics, verifying the physical stability of the MuJoCo engine.
\end{itemize}

% ================================================================
\section{Significance of the Findings}
\label{sec:significance}
% ================================================================

\begin{itemize}
    \item \textbf{Technical significance}: This project demonstrates the full implementation of multi-agent deep RL on a physically realistic continuous environment. The custom reward shaping wrapper is a concrete, reusable engineering contribution applicable to other sparse-reward MuJoCo tasks.
    
    \item \textbf{Practical significance}: The training infrastructure developed—including the Google Colab notebook with automatic Drive checkpointing—provides a practical, cost-free solution for training deep RL agents that would otherwise require dedicated local GPU clusters.
    
    \item \textbf{Scientific significance}: The results confirm theoretical predictions: PPO with parameter sharing and dense reward wrappers is highly effective for continuous multi-agent cooperative-competitive tasks.
\end{itemize}

% ================================================================
\section{Limitations}
\label{sec:limitations}
% ================================================================

\begin{itemize}
    \item \textbf{Training duration}: Free-tier Google Colab limits continuous sessions to approximately 3.8 hours, restricting the number of PPO timesteps achievable. Deep RL research typically requires tens of millions of steps; our budget is an order of magnitude smaller.
    
    \item \textbf{Policy Symmetry}: Parameter sharing assumes agent symmetry—all agents share the same policy network. This prevents the emergence of highly specialized asymmetric roles (e.g., goalkeeper vs. attacker). Asymmetric MARL approaches such as MAPPO~\citep{yu2022mappo} could address this.
    
    \item \textbf{Reward Coefficient Tuning}: The three reward coefficients ($-0.01$, $+0.10$, $+0.50$) were set by engineering judgment rather than a systematic hyperparameter sweep.
    
    \item \textbf{Evaluation sample size}: Only 10 evaluation episodes were run due to the computational budget, carrying some variance in the statistical conclusions.
\end{itemize}

% ================================================================
\section{Summary}
\label{sec:discussion-summary}
% ================================================================

This chapter provided an analytical interpretation of our multi-agent soccer results. PPO with parameter sharing and reward shaping produced a functional, highly competitive multi-agent soccer training pipeline. The limitations—primarily training duration and policy symmetry—are clearly understood. Overall, the results \textbf{support} the initial research objectives, demonstrating that continuous deep reinforcement learning algorithms, when configured with appropriate reward wrappers, can train robust cooperative multi-agent behavior in complex physical environments.